\documentclass[11pt]{article}

\usepackage[margin=1in]{geometry}
\usepackage{amsmath,amssymb,amsthm,bm}
\usepackage{array,booktabs,tabularx,longtable}
\usepackage{algorithm,algpseudocode}
\usepackage{enumitem}
\usepackage[dvipsnames]{xcolor}
\usepackage[colorlinks=true,linkcolor=MidnightBlue,citecolor=MidnightBlue,urlcolor=MidnightBlue]{hyperref}

% ---- notation ----------------------------------------------------------------
\newcommand{\R}{\mathbb{R}}
\newcommand{\E}{\mathbb{E}}
\newcommand{\Ls}{L_S}                    % training loss
\newcommand{\Ld}{L_{\mathcal{D}}}        % population loss
\newcommand{\lmax}{\lambda_{\max}}
\newcommand{\tr}{\operatorname{tr}}
\newcommand{\N}{\mathcal{N}}
\newcommand{\gp}{g_{\mathrm{adv}}}
\newcommand{\pp}{\,pp}                   % percentage points
\newcommand{\presenter}[1]{\noindent\textit{Presenter: #1}\par\medskip}

\newcolumntype{L}{>{\raggedright\arraybackslash}X}
\setlist{itemsep=2pt,topsep=4pt}
\renewcommand{\arraystretch}{1.15}
\theoremstyle{plain}
\newtheorem{theorem}{Theorem}
\newtheorem{lemma}{Lemma}

\title{Flat vs.\ Sharp Minima and Generalization\\[4pt]
\large E0 334 (Deep Learning for NLP), IISc --- Paper Presentation}
\author{Sai Aditya \and Sai Anoop \and Sai Vinay}
\date{27 and 29 October 2026}

\begin{document}
\maketitle
\tableofcontents
\newpage

% ============================================================================
\section{Overview}
% ============================================================================

\subsection{The question}

A modern network can reach almost zero training loss at many different weight settings, and these
settings can differ a lot in test accuracy. A long-standing explanation is the geometry of the loss
around the solution:

\begin{quote}
\textbf{Hypothesis (H).} \emph{Among weights with low training loss, those in \textbf{flat} regions (the loss stays low when
the weights are perturbed) generalize better than those in \textbf{sharp} regions (the loss rises quickly).}
\end{quote}

The papers we present attack (H) from three directions:
\begin{enumerate}
  \item \textbf{Build an optimizer that seeks flatness, then check accuracy}: Entropy-SGD, SAM and its variants.
  \item \textbf{Take the best such optimizer to language models}: SAM for T5 (Bahri et al.).
  \item \textbf{Test whether a flatness measure actually predicts generalization}: the maximum Hessian eigenvalue (Kaur et al.).
\end{enumerate}

\subsection{Papers and presenters}

\begin{center}\small
\begin{tabularx}{\textwidth}{@{}p{2.3cm}LLp{2.6cm}@{}}
\toprule
\textbf{Part / presenter} & \textbf{Papers} & \textbf{Notion of flatness} & \textbf{Domain} \\
\midrule
I.\ Entropy-SGD \newline \emph{Sai Anoop} &
Chaudhari et al., \emph{Entropy-SGD: Biasing Gradient Descent into Wide Valleys} (ICLR 2017) &
Local entropy: how much low-loss volume surrounds $w$ &
MNIST, CIFAR-10, PTB, char-level LM \\
\addlinespace
II.\ SAM family \newline \emph{Sai Aditya} &
Foret et al.\ (SAM, ICLR 2021); Bahri et al., \emph{SAM Improves Language Model Generalization} (ACL 2022);
GSAM (ICLR 2022); WSAM (KDD 2023); Eigen-SAM (2025); Duan et al., \emph{X-SAM} (2026) &
Worst-case loss in a ball of radius $\rho$; later variants target $\lmax$ directly &
GLUE, SuperGLUE, QA (T5); CIFAR, Fashion-MNIST \\
\addlinespace
III.\ Max.\ Hessian eigenvalue \newline \emph{Sai Vinay} &
Kaur, Cohen, Lipton, \emph{On the Maximum Hessian Eigenvalue and Generalization} (2023) &
$\lmax$, the top eigenvalue of the Hessian &
CIFAR-10, Fashion-MNIST \\
\bottomrule
\end{tabularx}
\end{center}

\subsection{Lecture plan (proposed)}

\begin{center}\small
\begin{tabular}{@{}llll@{}}
\toprule
\textbf{Lecture} & \textbf{Segment} & \textbf{Presenter} & \textbf{Time} \\
\midrule
1 & Overview and background (Sections 1--2)             & shared     & 15 min \\
1 & Part I: Entropy-SGD --- ideas, math, experiments     & Sai Anoop  & 40 min \\
1 & Part II(a): SAM and SAM for language models          & Sai Aditya & 35 min \\
\midrule
2 & Part II(b): GSAM, WSAM, Eigen-SAM, X-SAM             & Sai Aditya & 30 min \\
2 & Part III: the maximum Hessian eigenvalue             & Sai Vinay  & 35 min \\
2 & Synthesis and our contribution                       & all        & 20 min \\
2 & Questions                                            &            & 5 min \\
\bottomrule
\end{tabular}
\end{center}

% ============================================================================
\section{Background: loss landscapes and how to measure flatness}
\label{sec:background}
% ============================================================================

\subsection{Notation}

\begin{itemize}
  \item Weights $w \in \R^d$. Training set $S = \{(x_i,y_i)\}_{i=1}^n$ drawn from a distribution $\mathcal{D}$.
        Per-example loss $\ell(w;x,y)$ (cross-entropy).
  \item Training loss $\Ls(w) = \frac1n\sum_i \ell(w;x_i,y_i)$; population (test) loss $\Ld(w) = \E_{(x,y)\sim\mathcal{D}}\,\ell(w;x,y)$.
        The \textbf{generalization gap} is $\Ld(w) - \Ls(w)$. We write $L$ for $\Ls$ when the context is clear.
  \item Gradient $g = \nabla L(w)$; mini-batch gradient $\nabla L_B(w)$; learning rate $\eta$; batch size $B$.
  \item Hessian $H = \nabla^2 L(w) \in \R^{d\times d}$ with eigen-decomposition $H = \sum_{i=1}^d \lambda_i v_i v_i^\top$,
        $\lambda_1 = \lmax \ge \lambda_2 \ge \dots$, and unit eigenvectors $v_i$.
\end{itemize}

\subsection{Curvature is what ``flat'' and ``sharp'' mean}

Near $w$, a second-order Taylor expansion gives
\begin{equation}
  L(w+\delta) \approx L(w) + g^\top \delta + \tfrac12\,\delta^\top H \delta .
  \label{eq:taylor}
\end{equation}
At a minimum $g = 0$, so the loss increase under a perturbation is decided by $H$. Moving a distance $\rho$ along
eigenvector $v_i$ raises the loss by $\tfrac12 \rho^2 \lambda_i$. A minimum is \emph{sharp} along directions with large $\lambda_i$
and \emph{flat} along directions with $\lambda_i \approx 0$. Different flatness measures summarize the spectrum differently:

\begin{center}\small
\begin{tabularx}{\textwidth}{@{}p{3.6cm}p{5.4cm}Lp{2.2cm}@{}}
\toprule
\textbf{Measure} & \textbf{Definition} & \textbf{Meaning at a minimum} & \textbf{Used by} \\
\midrule
Top eigenvalue $\lmax$ & $\max_{\|u\|=1} u^\top H u$ & Loss increase in the \emph{worst} direction: $\max_{\|\delta\|\le\rho} L(w+\delta) - L(w) \approx \tfrac12\rho^2\lmax$ & SAM family, Kaur \\
Trace $\tr(H)$ & $\sum_i \lambda_i$ & Loss increase under \emph{random} noise: $\E_{\delta\sim\N(0,\sigma^2 I)}[L(w+\delta)] - L(w) \approx \tfrac12\sigma^2\tr(H)$ & X-SAM (spectrum) \\
Full spectrum & all $\lambda_i$ & How many directions are flat vs.\ sharp & Entropy-SGD, X-SAM \\
Local entropy & $\log\!\int e^{-L(w') - \frac{\gamma}{2}\|w-w'\|^2}\,dw'$ & Low-loss \emph{volume} around $w$ & Entropy-SGD \\
\bottomrule
\end{tabularx}
\end{center}

\subsection{Computing curvature without forming the Hessian}

For a model with $d \sim 10^8$ weights, $H$ has $d^2$ entries and cannot be stored. Everything above is computed
from \textbf{Hessian-vector products} (Pearlmutter, 1994):
\begin{equation}
  Hv = \nabla_w\big(\nabla_w L(w)^\top v\big),
\end{equation}
i.e.\ differentiate the scalar $g^\top v$ once more. This costs about two backward passes.
\begin{itemize}
  \item \textbf{Power iteration for $\lmax$:} start from a random unit $v$, repeat $v \leftarrow Hv / \|Hv\|$. The Rayleigh quotient
        $v^\top H v$ converges to $\lmax$ (assuming the top eigenvalue is also the largest in magnitude).
  \item \textbf{Hutchinson's estimator for the trace:} $\tr(H) = \E_z[z^\top H z]$ for random $z$ with i.i.d.\ $\pm1$ entries.
\end{itemize}

\subsection{Why flat minima might generalize}

\begin{itemize}
  \item \textbf{Robustness.} Training and test losses are slightly different functions of $w$. At a flat minimum,
        small shifts between the two surfaces change the loss little (Hochreiter and Schmidhuber, 1997).
  \item \textbf{Description length / Bayesian view.} A flat minimum can be specified with fewer bits of precision. Under a prior,
        a wide valley carries more probability mass than a narrow one of the same depth.
  \item \textbf{PAC-Bayes bounds.} The test loss of weights perturbed by noise is bounded by the perturbed training loss plus a complexity
        term. Flat minima keep the perturbed training loss low. This is the formal motivation for SAM (Part II).
  \item \textbf{But:} sharpness is not invariant to reparameterization. One can rescale a ReLU network so the function is unchanged but $\lmax$
        is arbitrarily large (Dinh et al., 2017). Part III takes this further with realistic training runs.
\end{itemize}

% ============================================================================
\section{Part I --- Entropy-SGD}
\label{sec:esgd}
% ============================================================================
\presenter{Sai Anoop}

Chaudhari, Choromanska, Soatto, LeCun, Baldassi, Borgs, Chayes, Sagun, Zecchina. \emph{Entropy-SGD: Biasing Gradient Descent into Wide Valleys}. ICLR 2017.

\subsection{Motivation: what the Hessian looks like at a good solution}

The authors compute the full Hessian spectrum at the end of training for small networks (Section~\ref{sec:esgd-setup}).
For a LeNet-style CNN on MNIST:
\begin{itemize}
  \item about 94\% of the eigenvalues have $|\lambda| < 10^{-4}$, so almost every direction is flat;
  \item a long positive tail of a few large eigenvalues (up to $\approx 40$);
  \item a very short negative tail (most negative $\approx -0.46$): the descent directions the optimizer missed are nearly flat too.
\end{itemize}
The same shape appears for an MLP, an LSTM and a larger CNN. So good solutions found by SGD lie in \emph{wide valleys}.
The question is: \emph{can we modify SGD to seek such regions on purpose?}

\subsection{From the Gibbs distribution to local entropy}

The Gibbs distribution of a loss $f$ at inverse temperature $\beta$ is
\begin{equation}
  P(w;\beta) = Z_\beta^{-1}\, e^{-\beta f(w)} .
\end{equation}
As $\beta\to\infty$ it concentrates on the global minimum, so minimizing $f$ is the same as finding its mode. To focus on
\emph{wide} regions instead, add a coupling that keeps the distribution near a reference point $w$ (with $\beta=1$):
\begin{equation}
  P(w'; w, \gamma) \propto \exp\!\Big(-f(w') - \frac{\gamma}{2}\|w - w'\|^2\Big).
  \label{eq:local-gibbs}
\end{equation}

\paragraph{Definition (local entropy).}
\begin{equation}
  F(w,\gamma) = \log \int_{\R^d} \exp\!\Big(-f(w') - \frac{\gamma}{2}\|w - w'\|^2\Big)\, dw' .
  \label{eq:local-entropy}
\end{equation}
$F$ is the log-partition function of~\eqref{eq:local-gibbs}. It is large when many $w'$ near $w$ have low loss, so it measures
the \emph{depth and width} of the valley at once. Entropy-SGD solves
\begin{equation}
  \min_w\; -F(w,\gamma).
\end{equation}

\paragraph{The scope $\gamma$.} Large $\gamma$ confines $w'$ to a tiny neighbourhood and $-F(w,\gamma) \to f(w) + \text{const}$
(the original problem). Small $\gamma$ averages over a wide region, so the landscape of $-F$ is very smooth. In the paper's 1-D
cartoon, a shallow but wide minimum becomes the \emph{global} minimum of $-F$, even though a deeper sharp minimum exists.

\paragraph{Why not classical entropy or plain smoothing?} Classical entropy would favour any wide region, including wide plateaus
of high loss. Convolving $f$ with a kernel can create a spurious minimum between two sharp valleys. Local entropy weights each
region by $e^{-f}$, so it rewards regions that are \emph{both} low and wide.

\subsection{The gradient of local entropy}

Differentiate~\eqref{eq:local-entropy} under the integral sign. With $Z(w) = \int e^{-f(w') - \frac{\gamma}{2}\|w-w'\|^2} dw'$,
\begin{equation}
  \nabla_w F = \frac{1}{Z(w)} \int -\gamma (w - w')\, e^{-f(w') - \frac{\gamma}{2}\|w-w'\|^2}\, dw'
            = -\gamma\big(w - \langle w' \rangle\big),
  \qquad \langle w' \rangle = \E_{P(w';w,\gamma)}[w'] .
  \label{eq:entropy-grad}
\end{equation}
So the descent direction for $-F$ is $\gamma(w - \langle w'\rangle)$: \emph{move $w$ towards the mean of the local Gibbs distribution}.
When two nearby minima have similar loss, $\langle w'\rangle$ lies closer to the wider one, so the step points towards wide valleys.

\paragraph{Worked example (ours): a quadratic valley.} Take $f(w') = \frac{a}{2} w'^2$ in 1-D. Completing the square in~\eqref{eq:local-entropy},
\begin{equation}
  -F(w,\gamma) = \frac12\,\frac{a\gamma}{a+\gamma}\, w^2 + \text{const},
  \qquad \langle w'\rangle = \frac{\gamma}{a+\gamma}\, w .
\end{equation}
The effective curvature $\frac{a\gamma}{a+\gamma}$ is always below $\min(a,\gamma)$. Sharp directions ($a \gg \gamma$) are \emph{capped} at about $\gamma$,
while flat directions ($a \ll \gamma$) are almost unchanged. Local entropy therefore removes sharp curvature, and $\gamma$ sets the cap.

\subsection{Estimating $\langle w'\rangle$ with Langevin dynamics}

The expectation in~\eqref{eq:entropy-grad} is over a high-dimensional distribution, so it is estimated by MCMC.
Stochastic gradient Langevin dynamics (SGLD; Welling and Teh, 2011) samples from $p(w') \propto e^{-U(w')}$ with noisy gradient steps:
\begin{equation}
  w' \leftarrow w' - \eta' \nabla U(w') + \sqrt{\eta'}\,\varepsilon\, \xi, \qquad \xi\sim\N(0,I).
\end{equation}
Here $U(w') = f(w') + \frac{\gamma}{2}\|w-w'\|^2$, so $\nabla U(w') = \nabla f(w') - \gamma(w - w')$. This is SGD on $f$ plus a spring pulling
$w'$ back to $w$, plus noise. Averaging the SGLD iterates gives $\mu \approx \langle w'\rangle$.

\begin{algorithm}[h]
\caption{Entropy-SGD (one outer step)}
\begin{algorithmic}[1]
\Require weights $w$; scope $\gamma$; outer lr $\eta$; SGLD step $\eta'$; thermal noise $\varepsilon$; SGLD steps $L$; averaging weight $\alpha$
\State $w' \gets w$, \; $\mu \gets w$
\For{$\ell = 1, \dots, L$}
  \State sample a mini-batch $B_\ell$
  \State $dw' \gets \nabla L_{B_\ell}(w') - \gamma\,(w - w')$
  \State $w' \gets w' - \eta'\, dw' + \sqrt{\eta'}\,\varepsilon\,\N(0,I)$
  \State $\mu \gets (1-\alpha)\,\mu + \alpha\, w'$
\EndFor
\State $w \gets w - \eta\,\gamma\,(w - \mu)$ \Comment{in practice $w \gets w - \eta\,(w-\mu)$, see scoping}
\end{algorithmic}
\end{algorithm}

\paragraph{Practical details.}
\begin{itemize}
  \item ``SGD inside SGD'': each outer update costs $L$ gradient evaluations ($L \in [5, 20]$). Comparisons use \emph{effective epochs} $=$ epochs $\times L$.
  \item $\alpha = 0.75$ puts more weight on later samples, which acts like burn-in. $\eta' \in [0.1, 1]$ is kept fixed (annealing it hurt).
        $\varepsilon \in [10^{-4}, 10^{-3}]$.
  \item Momentum (Nesterov) or Adam can be used in either loop. Entropy-Adam feeds $w - \mu$ to Adam as the gradient.
  \item \textbf{Scoping:} $\gamma$ grows during training, $\gamma_t = \gamma_0 (1+\gamma_1)^t$, so the search starts coarse (smooth landscape) and becomes finer.
        Because a growing $\gamma$ would interfere with the learning-rate schedule, the factor $\gamma$ is dropped from the outer step.
  \item Tuning heuristic: pick $\gamma$ so that $\|\gamma(w-\mu)\|$ is comparable to the plain SGD gradient.
\end{itemize}

\subsection{Theory: a smoother objective and a stability bound}

Assume $f$ is $\alpha$-Lipschitz and $\beta$-smooth, and (for the analysis) that no Hessian eigenvalue lies in $[-2\gamma - c,\, c]$ for some $c>0$.

\begin{lemma}[Lemma 2]
$-F(\cdot,\gamma)$ is $\frac{\alpha}{1+\gamma^{-1}c}$-Lipschitz and $\frac{\beta}{1+\gamma^{-1}c}$-smooth, i.e.\ smoother than $f$.
\end{lemma}

\begin{theorem}[Hardt et al., 2016]
Uniform stability bounds generalization: if changing one training example changes the loss of the output by at most $\epsilon$, the expected
generalization gap is at most $\epsilon$. For SGD on an $\alpha$-Lipschitz, $\beta$-smooth loss with step $\eta_t \le 1/t$ for $T$ steps on $N$ samples,
$\epsilon \lesssim \frac1N\, \alpha^{1/(1+\beta)}\, T^{1 - 1/(1+\beta)}$.
\end{theorem}

Substituting the smaller constants of Lemma~2 gives a smaller stability bound for Entropy-SGD than for SGD (their Eq.~9).
\textbf{Caveat (the authors' Remark 4):} the eigenvalue-gap assumption is ``admittedly unrealistic'', since the motivating experiment shows
most eigenvalues are near zero. The bound is an intuition, not a guarantee.

\subsection{Experimental setup and results}
\label{sec:esgd-setup}

\paragraph{(A) Hessian at the solutions of standard optimizers (observational).}
\begin{center}\small
\begin{tabularx}{\textwidth}{@{}p{4.0cm}p{2.2cm}p{2.6cm}L@{}}
\toprule
\textbf{Network} & \textbf{Data} & \textbf{Trained with} & \textbf{Measured} \\
\midrule
small-LeNet (47,658 weights) & MNIST & Adam (2.4\% error) & exact Hessian spectrum \\
small-mnistfc (50,890 weights) & MNIST & SGD + momentum (2.5\%) & exact Hessian spectrum \\
char-LSTM (32,640 weights) & small text & Adam & exact Hessian spectrum \\
All-CNN-BN ($\approx$1.6M weights) & CIFAR-10 & Adam (11.2\%) & diagonal Fisher (Hessian proxy) \\
\bottomrule
\end{tabularx}
\end{center}

\paragraph{(B) Entropy-SGD vs.\ standard optimizers.}
\begin{center}\small
\begin{tabularx}{\textwidth}{@{}p{2.3cm}p{2.9cm}p{3.0cm}L@{}}
\toprule
\textbf{Data} & \textbf{Model} & \textbf{Baseline} & \textbf{Entropy-SGD} \\
\midrule
MNIST & mnistfc ($2\times1024$), LeNet + BN & Adam, 100 epochs & $L{=}20$, $\eta'{=}0.1$, $\varepsilon{=}10^{-3}$, $\gamma_0{=}10^{-4}$, $\times1.001$/step \\
CIFAR-10 (no augmentation) & All-CNN-C + BN & SGD + Nesterov, 200 epochs & $L{=}20$, $\eta'{=}0.1$, $\varepsilon{=}10^{-4}$, $\gamma_0{=}0.03$, $\times1.001$/step \\
PTB (word LM) & 2-layer LSTM, 66M weights & SGD (Zaremba et al.) & $L{=}5$, fixed $\gamma{=}10^{-3}$, $\eta'{=}1$, momentum 0.5 \\
\emph{War and Peace} (char LM) & 2-layer LSTM ($128$ units) & Adam, 50 epochs & Entropy-Adam, $L{=}5$, fixed $\gamma{=}0.01$, $\beta_1{=}0.5$ \\
\bottomrule
\end{tabularx}
\end{center}

\begin{center}\small
\begin{tabular}{@{}llll@{}}
\toprule
\textbf{Model} & \textbf{Metric} & \textbf{Entropy-SGD (eff.\ epochs)} & \textbf{Baseline (epochs)} \\
\midrule
mnistfc     & test error (\%)  & $1.37 \pm 0.03$ (120)  & $1.39 \pm 0.03$ (66) \\
LeNet       & test error (\%)  & $0.50 \pm 0.01$ (80)   & $0.51 \pm 0.01$ (100) \\
All-CNN-BN  & val.\ error (\%) & $7.81 \pm 0.09$ (160)  & $7.71 \pm 0.19$ (180) \\
PTB-LSTM    & test perplexity  & $77.66 \pm 0.17$ (25)  & $78.6 \pm 0.26$ (55) \\
char-LSTM   & test perplexity  & $1.217 \pm 0.005$ (25) & $1.226 \pm 0.01$ (40) \\
\bottomrule
\end{tabular}
\end{center}

Control: plain SGLD, without the local-entropy coupling, is much worse (0.63\% on LeNet, 9.89\% on All-CNN-BN).

\subsection{Critical reading}
\begin{itemize}
  \item On CNNs, Entropy-SGD \emph{matches} the baselines (slightly worse on CIFAR-10). The real wins are on LSTMs: better perplexity in about half the effective epochs.
  \item Experiment (A) is correlational. It never compares a flat and a sharp solution on the same task, and the flatness of Entropy-SGD's own solutions is not measured.
  \item The generalization theory rests on an assumption the paper's own data contradicts.
\end{itemize}

% ============================================================================
\section{Part II --- Sharpness-Aware Minimization and its variants}
\label{sec:sam}
% ============================================================================
\presenter{Sai Aditya}

\subsection{SAM (Foret et al., ICLR 2021)}

\paragraph{From a generalization bound to an objective.} A PAC-Bayes argument gives, with high probability over the training set,
\begin{equation}
  \Ld(w) \;\le\; \max_{\|\epsilon\|_2\le\rho} \Ls(w+\epsilon) \;+\; h\!\left(\|w\|_2^2/\rho^2\right),
\end{equation}
with $h$ increasing. The first term splits as
\begin{equation}
  \max_{\|\epsilon\|\le\rho} \Ls(w+\epsilon) = \underbrace{\Big[\max_{\|\epsilon\|\le\rho} \Ls(w+\epsilon) - \Ls(w)\Big]}_{\text{sharpness}} + \Ls(w),
\end{equation}
so minimizing it lowers the training loss \emph{and} the sharpness. SAM solves
\begin{equation}
  \min_w\; L^{\mathrm{SAM}}(w) + \lambda\|w\|_2^2, \qquad L^{\mathrm{SAM}}(w) = \max_{\|\epsilon\|_2\le\rho} \Ls(w+\epsilon).
\end{equation}

\paragraph{Solving the inner maximization.} Linearize, $L(w+\epsilon) \approx L(w) + \epsilon^\top g$. Maximizing a linear function over a ball gives
\begin{equation}
  \hat\epsilon(w) = \arg\max_{\|\epsilon\|_2\le\rho} \epsilon^\top g = \rho\,\frac{g}{\|g\|_2}
\end{equation}
(for a general $\ell_p$ ball the solution is the dual-norm direction $\rho\,\mathrm{sign}(g)|g|^{q-1}/\|g\|_q^{q-1}$, with $1/p + 1/q = 1$).

\paragraph{The gradient.} Differentiating $L(w + \hat\epsilon(w))$ gives
\begin{equation}
  \nabla_w L^{\mathrm{SAM}}(w) \approx \nabla L(w)\big|_{w+\hat\epsilon} + \Big(\frac{d\hat\epsilon}{dw}\Big)^{\!\top} \nabla L(w)\big|_{w+\hat\epsilon}.
\end{equation}
The second term contains the Hessian and is dropped. So SAM uses \textbf{the gradient at the adversarial point, applied at the original point}.

\begin{algorithm}[h]
\caption{SAM (one step, any base optimizer)}
\begin{algorithmic}[1]
\State $g \gets \nabla L_B(w)$ \Comment{1st forward/backward}
\State $w_{\mathrm{adv}} \gets w + \rho\, g/\|g\|$ \Comment{ascent step}
\State $\gp \gets \nabla L_B(w_{\mathrm{adv}})$ \Comment{2nd forward/backward}
\State $w \gets \mathrm{BaseOptimizer}(w, \gp)$ \Comment{e.g.\ $w - \eta\,\gp$}
\end{algorithmic}
\end{algorithm}

\paragraph{Link to $\lmax$.} At a minimum ($g=0$), by~\eqref{eq:taylor} the sharpness term is $\max_{\|\epsilon\|\le\rho} \frac12\epsilon^\top H\epsilon = \frac12\rho^2\lmax$.
So the \emph{idealized} SAM objective penalizes $\lmax$. The first-order algorithm, however, only sees $g$, which vanishes at a minimum.
This gap motivates every variant below.

\paragraph{Two practical knobs.}
\begin{itemize}
  \item \textbf{$m$-sharpness:} split the batch into $m$ pieces and compute a separate $\hat\epsilon$ for each; small pieces often generalize better.
  \item \textbf{ASAM} (Kwon et al., 2021): sharpness depends on weight scale, so measure the ball relative to $|w|$:
        $\hat\epsilon = \rho\, T_w^2 g / \|T_w g\|$ with $T_w = \mathrm{diag}(|w|)$.
\end{itemize}

\subsection{SAM for language models (Bahri, Mobahi, Tay; ACL 2022)}

\paragraph{Idea.} SAM had been shown on images only. Fine-tune public T5 checkpoints with and without SAM on standard NLP benchmarks.

\paragraph{Efficient SAM.} The ascent gradient is computed on a micro-batch $M$ of only $\frac14$ of the batch (after Brock et al., 2021):
\begin{equation}
  w_{\mathrm{adv}} = w + \rho\, \frac{\nabla L_M(w)}{\|\nabla L_M(w)\|}, \qquad w \leftarrow \mathrm{opt}\big(w, \nabla L_B(w_{\mathrm{adv}})\big),
\end{equation}
which makes SAM only about 25\% slower than standard training.

\paragraph{Setup.}
\begin{itemize}
  \item \textbf{Models:} T5.1.1 Small (77M), Base (250M), Large (800M), XL (3B); mT5 for multilingual QA.
  \item \textbf{Optimizer:} Adafactor, lr $10^{-3}$, batch 128; SAM with $\rho = 0.05$ (Small) or $0.15$ (others), tuned once on SuperGLUE.
  \item \textbf{Protocol:} one run per setting; best checkpoint chosen per task and metric.
\end{itemize}

\begin{center}\small
\begin{tabularx}{\textwidth}{@{}p{4.2cm}p{3.0cm}L@{}}
\toprule
\textbf{Benchmark} & \textbf{Training} & \textbf{Result (no SAM $\to$ SAM)} \\
\midrule
SuperGLUE (dev) & task mixture, 250k steps & improves at all sizes; Base $75.3\to78.5$, XL $87.2\to89.1$ \\
GLUE (dev) & task mixture, 250k steps & improves at all sizes; Large $86.2\to88.3$ \\
Closed-book QA: Natural Questions, WebQuestions, TriviaQA (test) & mixture or single task, 20k steps & improves for every task, metric and size; Base EM on NQ $18.1\to20.6$ \\
TyDiQA-GoldP (mT5, 9 languages) & 20k steps, $\rho \in \{0.02, 0.05\}$ & modest: $+0.3$ to $+0.9$ average F1 \\
Low data: train sets subsampled to 2--80\% & SuperGLUE, CBQA & gains persist; 5\% SuperGLUE Base $52.9\to56.7$; largest relative gains near 20\% \\
Ablations (SuperGLUE Base) & $\rho$, micro-batch $a$, $m$ & every $\rho$ beats no SAM, $\rho{=}0.15$ best; $a = 32$ of 128 suffices; $m>1$ does not help \\
\bottomrule
\end{tabularx}
\end{center}

\paragraph{Critical reading.} Strong evidence that SAM helps NLP fine-tuning, especially with little data. But \textbf{flatness is never measured},
so the paper shows SAM \emph{helps}, not that it helps \emph{because} the solution is flatter. Results are single-seed, and a few tasks get worse
(e.g.\ GLUE RTE for Small $74.0\to70.4$).

\subsection{Why plain SAM is not enough: the surrogate gap (GSAM)}

Zhuang et al.\ (ICLR 2022) point out that a low perturbed loss $f_p(w) = \max_{\|\delta\|\le\rho} f(w+\delta)$ can occur at a sharp minimum if that minimum is deep enough.
So for two minima, $f_p(w_1) \le f_p(w_2)$ does \emph{not} imply $\lmax(w_1) \le \lmax(w_2)$. A quantity that does track sharpness is the \textbf{surrogate gap}
\begin{equation}
  h(w) = f_p(w) - f(w) \approx f(w_{\mathrm{adv}}) - f(w) \;\ge\; 0, \qquad h(w^*) \approx \tfrac12\rho^2\lmax(w^*) \text{ at a minimum } w^*.
\end{equation}

\paragraph{GSAM.} Minimize $f_p$ and $h$ together, without letting the second objective hurt the first. Decompose the clean gradient $g$ relative to the SAM gradient $\gp$:
\begin{equation}
  g = g_\parallel + g_\perp, \qquad g_\parallel = \frac{\langle g, \gp\rangle}{\|\gp\|^2}\,\gp .
\end{equation}
Since $\nabla h = \gp - g$, its component orthogonal to $\gp$ is $-g_\perp$. Moving along $+g_\perp$ therefore lowers $h$ and leaves $f_p$ unchanged to first order:
\begin{equation}
  w \leftarrow w - \eta\,(\gp - \alpha\, g_\perp).
\end{equation}
GSAM also ties the radius to the learning-rate schedule: $\rho_t = \rho_{\min} + (\rho_{\max}-\rho_{\min})\frac{\eta_t - \eta_{\min}}{\eta_{\max}-\eta_{\min}}$.

\subsection{WSAM: sharpness as a weighted regularizer (Yue et al., KDD 2023)}

Make the trade-off between loss and sharpness explicit:
\begin{equation}
  L_{\mathrm{WSAM}}(w) = L(w) + \frac{\gamma}{1-\gamma}\big(L_{\mathrm{SAM}}(w) - L(w)\big)
                      = \frac{1-2\gamma}{1-\gamma}L(w) + \frac{\gamma}{1-\gamma}L_{\mathrm{SAM}}(w), \qquad \gamma\in[0,1).
\end{equation}
$\gamma = 0$ is plain training, $\gamma = \frac12$ is SAM, and $\gamma > \frac12$ weights sharpness more than SAM does. Like decoupled weight decay in AdamW,
the base optimizer sees only the clean gradient $\tilde g$, and the sharpness gradient is applied separately:
\begin{equation}
  w \leftarrow \mathrm{BaseStep}(w, \tilde g) - \eta\,\frac{\gamma}{1-\gamma}\,(\gp - \tilde g).
\end{equation}

\subsection{Eigen-SAM: why SAM only \emph{implicitly} controls $\lmax$ (Luo et al., 2025)}

\paragraph{Continuous-time view.} SGD is well approximated by the SDE $dW_t = -\nabla f(W_t)\,dt + \sqrt{\eta}\,\Sigma^{1/2} dB_t$.
Expanding SAM's perturbed gradient to \emph{third} order gives an SDE whose drift is the gradient of
\begin{equation}
  \tilde f_{\mathrm{SAM}}(w) = f(w) + \rho\, \E\|\nabla f_\gamma(w)\| + \frac{\rho^2}{2}\, \E\!\left[\frac{\nabla f_\gamma^\top\, \nabla^2 f_\gamma(w)\, \nabla f_\gamma}{\|\nabla f_\gamma\|^2}\right],
\end{equation}
where $f_\gamma$ is the mini-batch loss. The last term is the curvature \emph{along the gradient direction}. If the perturbation is aligned with the top eigenvector $v_1$, this term becomes $\frac{\rho^2}{2}\lmax$,
and SAM implicitly minimizes $\lmax$. Measured in practice, this alignment is weak, so the implicit regularization is weak.
A PAC-Bayes bound in the same paper bounds the test loss by the training loss plus a term proportional to $\lmax$ (plus higher-order and complexity terms).

\paragraph{Eigen-SAM.} Every $p$ steps, estimate $v$ by $q$ power iterations. Then tilt the perturbation towards $v$:
\begin{equation}
  \epsilon = \rho\left(\frac{g}{\|g\|} + \alpha\,\mathrm{sign}\langle g, v\rangle\, v_\perp\right), \qquad v_\perp = v - \Big\langle v, \tfrac{g}{\|g\|}\Big\rangle \tfrac{g}{\|g\|},
\end{equation}
and update with $\nabla f(w + \epsilon)$. Default values are $\alpha = 0.2$, $p = 100$, $q = 5$; $\alpha = 0$ gives SAM.

\subsection{X-SAM: correct the update, not the perturbation (Duan et al., 2026)}

\paragraph{Diagnostic.} The authors train a 6-layer CNN on CIFAR-10 with SAM and record the angle between SAM's update and $v_1$.
In about 45\% of steps the angle is $80^\circ$--$100^\circ$, i.e.\ nearly orthogonal.

\paragraph{Why the angle matters.} The quadratic change of the loss after a step $\Delta w = -\eta g$ has a contribution from the top direction of
\begin{equation}
  \lambda_1\, (\Delta w^\top v_1)^2 = \lambda_1\, \eta^2 \langle g, v_1\rangle^2 .
\end{equation}
\begin{itemize}
  \item If $g \perp v_1$, the step barely interacts with the sharpest direction, so $\lmax$ is hardly reduced. Their Theorem 5.3 makes this precise:
        $|\Delta\lambda_1| \le \eta\,\varepsilon\,\|g\| + O(\eta^2\|g\|^2/\text{gap})$, where $\varepsilon$ bounds the part of $\nabla\lambda_1$ orthogonal to $v_1$.
  \item If $g$ is partly aligned with $v_1$ (angles of $50^\circ$--$80^\circ$), the step moves along the sharpest direction, and the authors argue that this keeps the weights in sharp regions.
\end{itemize}

\paragraph{X-SAM update.} Normalize the SAM gradient and remove a fraction $\alpha$ of its component along $v$:
\begin{equation}
  \hat h = \frac{\gp}{\|\gp\|}, \qquad
  g_{\mathrm{X\text{-}SAM}} = \hat h - \alpha\,\mathrm{sign}\langle \gp, v\rangle\, \hat h_\parallel
  \;=\; (I - \alpha\, v v^\top)\,\hat h \quad (\text{projector form}),
\end{equation}
then $w \leftarrow w - \eta\, g_{\mathrm{X\text{-}SAM}}$. The eigenvector $v$ is re-estimated every $p$ steps, as in Eigen-SAM.

\paragraph{Theory.} With $P = I - \alpha vv^\top$ and $\alpha\in[0,2]$, $\|P\|_{\mathrm{op}}\le 1$. X-SAM converges at the $O(1/\sqrt T)$ rate of SGD
for smooth non-convex losses, up to a term from the perturbation that grows with $\rho$ (Theorem 5.1). A matrix-perturbation bound shows that the corrected step
can lower $\lmax$ (Theorem 5.2).

\paragraph{Our implementation note.} Since $\hat h_\parallel = \langle\hat h, v\rangle v$ does not depend on the sign of $v$, the sign factor in their Algorithm 2 can
\emph{add} to the sharp component for some sign combinations. Our code uses the projector form $(I - \alpha vv^\top)\hat h$, which always shrinks it.

\paragraph{Setup.} CIFAR-10, CIFAR-100 and Fashion-MNIST with standard augmentation. Models: ResNet-18, ResNet-50, WideResNet-28-10 and AlexNet.
Optimizer: SGD with lr 0.1 (cosine schedule), momentum 0.9, batch 256, 200 epochs. $\rho = 0.05$ (0.1 for CIFAR-100) and $\alpha = 0.2$. Three seeds per setting.

\begin{center}\small
\begin{tabular}{@{}lccccc@{}}
\toprule
\textbf{Avg.\ top-1 accuracy (\%)} & \textbf{SAM} & \textbf{WSAM} & \textbf{GSAM} & \textbf{Eigen-SAM} & \textbf{X-SAM} \\
\midrule
CIFAR-10 (4 models)  & 92.41 & 92.25 & 92.31 & 92.57 & \textbf{93.68} \\
CIFAR-100 (4 models) & 71.06 & 71.95 & 70.86 & 71.08 & \textbf{73.81} \\
Fashion-MNIST        & 94.90 & 95.43$^\ast$ & 95.16 & 95.24$^\ast$ & 95.24 \\
\bottomrule
\end{tabular}\\[2pt]
{\footnotesize $^\ast$Did not converge on AlexNet, so the average covers 3 models.}
\end{center}

They also measure flatness directly. For ResNet-18, X-SAM's $\lmax$ falls faster and oscillates less during training. Its Hessian spectrum has a smaller $\lmax$, a smaller trace, and more mass near zero than SAM's.

\paragraph{Critical reading.}
\begin{itemize}
  \item This is the one paper in Part II whose method targets $\lmax$, measures $\lmax$, \emph{and} reports accuracy.
  \item Much of the average gain comes from AlexNet ($+8.6\pp$ over SAM on CIFAR-100). On the ResNets the gains are $0.1$--$2\pp$, often within one standard deviation.
  \item Some numbers are inconsistent. The WSAM Fashion-MNIST average (95.43) does not match its row (95.10), and the headline gains differ between the abstract and Section 6.2.
  \item Vision only. The authors list transformers as future work.
\end{itemize}

\subsection{The SAM family at a glance}

\begin{center}\small
\begin{tabularx}{\textwidth}{@{}p{2.2cm}p{4.6cm}Lp{2.2cm}@{}}
\toprule
\textbf{Method} & \textbf{Perturbation} & \textbf{Update direction} & \textbf{Targets} \\
\midrule
SAM       & $\rho\, g/\|g\|$ & $\gp$ & worst-case loss \\
ASAM      & $\rho\, T_w^2 g/\|T_w g\|$ & $\gp$ & scale-invariant sharpness \\
GSAM      & $\rho_t\, g/\|g\|$ & $\gp - \alpha g_\perp$ & surrogate gap $h \approx \frac12\rho^2\lmax$ \\
WSAM      & $\rho\, g/\|g\|$ & $\tilde g$, then a separate $\frac{\gamma}{1-\gamma}(\gp - \tilde g)$ step & weighted sharpness \\
Eigen-SAM & $\rho\,(g/\|g\| + \alpha\,\mathrm{sign}\langle g,v\rangle v_\perp)$ & $\gp$ & $\lmax$, implicitly \\
X-SAM     & $\rho\, g/\|g\|$ & $(I - \alpha vv^\top)\,\gp/\|\gp\|$ & $\lmax$, explicitly \\
\bottomrule
\end{tabularx}
\end{center}

% ============================================================================
\section{Part III --- Is $\lmax$ really linked to generalization?}
\label{sec:kaur}
% ============================================================================
\presenter{Sai Vinay}

Kaur, Cohen, Lipton. \emph{On the Maximum Hessian Eigenvalue and Generalization}. arXiv:2206.10654 (v3, 2023).

\subsection{$\lmax$ as a sharpness measure}

$\lmax = \max_{\|u\|=1} u^\top H u$ is the curvature in the worst direction. It connects to the two most common definitions of sharpness:
\begin{itemize}
  \item Keskar et al.'s sharpness at a minimum is a scaled $\lmax$ if the loss is close to its quadratic approximation.
  \item SAM's idealized objective is bounded using $\lmax$. By~\eqref{eq:taylor} and $\epsilon^\top H\epsilon \le \|\epsilon\|^2\lmax$,
  \begin{equation}
    \max_{\|\epsilon\|\le\rho} L(w+\epsilon) \;\lesssim\; L(w) + \rho\,\|\nabla L(w)\| + \tfrac12\rho^2\,\lmax .
  \end{equation}
\end{itemize}

\subsection{Earlier doubt: sharpness is not invariant to reparameterization}

For a ReLU network, $\mathrm{relu}(cz) = c\,\mathrm{relu}(z)$ for $c>0$. So scaling one layer by $c$ and the next by $1/c$ leaves the function unchanged:
$L(c\theta_1, \theta_2/c) = L(\theta_1,\theta_2)$. Differentiating this identity twice with respect to the second layer shows that the Hessian block of the second layer at the
rescaled point is $c^2$ times the original. As $c\to\infty$, $\lmax\to\infty$ while the predictions stay identical (Dinh et al., 2017).
\textbf{Objection:} SGD never produces such rescaled weights. Kaur et al.\ ask whether $\lmax$ fails even for solutions SGD \emph{does} reach.

\subsection{What controls $\lmax$ in practice}

\paragraph{The edge of stability.} Gradient descent on a quadratic with curvature $\lambda$ updates the error as $\delta_{t+1} = (1-\eta\lambda)\,\delta_t$.
This is stable only if $|1-\eta\lambda| < 1$, i.e.\ $\lambda < 2/\eta$. Full-batch GD on neural networks first sharpens until $\lmax \approx 2/\eta$
and then hovers there (Cohen et al., 2021). So for large batches \textbf{the learning rate alone sets $\lmax$}, whatever the effect on test accuracy.

\paragraph{The linear scaling rule.} In the small-batch regime, SGD's behaviour (and its generalization) depends mostly on the ratio of batch size to learning rate.
Doubling both leaves accuracy unchanged, but the larger learning rate still pushes $\lmax$ down.

\paragraph{Other levers.} SAM's $\rho$, dropout noise, and batch normalization. For BN networks at small $\eta$, $\lmax$ is measured as the maximum over 8 points
between two successive iterates, because it varies a lot along the path.

\subsection{The logic of the paper}

If small $\lmax$ \emph{caused} good generalization, two things should hold:
\begin{enumerate}
  \item interventions that lower $\lmax$ should improve test accuracy;
  \item solutions with similar $\lmax$ should have similar test accuracy.
\end{enumerate}
The paper gives five counterexamples from ordinary training.

\subsection{Experimental setup and results}

\paragraph{Protocol.} Minibatch SGD with a constant learning rate and no schedule. Training stops at 99\% training accuracy. $\lmax$ is computed on a subset of the training set,
with dropout off. The same initialization seeds are used across settings, and results are averaged over 3--5 seeds.

\begin{center}\small
\begin{tabularx}{\textwidth}{@{}p{2.8cm}p{3.0cm}p{3.4cm}L@{}}
\toprule
\textbf{Intervention} & \textbf{Data / model} & \textbf{Varied} & \textbf{Finding} \\
\midrule
1.\ Larger learning rate & CIFAR-10, VGG-11 without BN & lr $10^{-4}$--$10^{-1}$; batch 100 to 50,000 (full) &
Larger lr always lowers $\lmax$ (for large batches $\lmax \approx 2/\eta$). Accuracy improves only for small batches. Runs with equal $\lmax$ can differ a lot in accuracy \\
\addlinespace
2.\ Linear scaling & CIFAR-10, VGG-11 (and Fixup ResNet-32) & batch 100/200/500 at a fixed batch/lr ratio &
Same accuracy, different $\lmax$ \\
\addlinespace
3.\ SAM & CIFAR-10, VGG-11, Fixup ResNet-32 & $\rho\in[0,0.1]$; small vs.\ large batch &
Higher $\rho$ always lowers $\lmax$, but accuracy rises only for small batches. With BN, SAM helps at both batch sizes \\
\addlinespace
4.\ Dropout & Fashion-MNIST, 2-hidden-layer MLP & $p \in [0, 0.7]$ &
Higher $p$ always lowers $\lmax$. Accuracy improves up to $p \approx 0.2$, then degrades \\
\addlinespace
5.\ Batch norm & CIFAR-10, VGG-11 with/without BN & lr sweep, both batch regimes &
$\lmax$ is about the same at a fixed lr, yet accuracy differs \\
\bottomrule
\end{tabularx}
\end{center}

\subsection{Critical reading}
\begin{itemize}
  \item Conclusion: $\lmax$ is not a reliable predictor of generalization, even for realistic SGD solutions.
  \item Not a claim that SAM or large learning rates are useless. Both still help with small batches. The claim is that $\lmax$ does not \emph{explain} why.
  \item Scope: one measure ($\lmax$; the trace is untested), small vision models, constant learning rates.
\end{itemize}

% ============================================================================
\section{Synthesis}
\label{sec:synthesis}
% ============================================================================

\begin{center}\small
\begin{tabularx}{\textwidth}{@{}p{3.8cm}LL@{}}
\toprule
\textbf{Question} & \textbf{Evidence for (H)} & \textbf{Evidence against / limits} \\
\midrule
Are good solutions flat? & Entropy-SGD: $\approx$95\% near-zero eigenvalues at SGD/Adam solutions & Correlational; no flat-vs-sharp comparison \\
\addlinespace
Does seeking flatness improve accuracy? & Entropy-SGD (LSTMs), SAM on T5 (GLUE, SuperGLUE, QA), X-SAM (CIFAR) & Entropy-SGD only matches baselines on CNNs; Bahri et al.\ never verify that flatness changed \\
\addlinespace
Does the flatness measure track accuracy? & X-SAM: lower $\lmax$ goes with higher accuracy across SAM variants & Kaur et al.: five interventions where $\lmax$ and accuracy move independently \\
\bottomrule
\end{tabularx}
\end{center}

\begin{enumerate}
  \item The papers make two different claims. \emph{Flat-seeking optimizers improve accuracy} is well supported.
        \emph{Flatness, measured as $\lmax$, is the reason} is contested.
  \item ``Flatness'' is several different quantities (local entropy, worst-case loss, $\lmax$, trace), and they can disagree.
  \item \textbf{None of these papers measures flatness for language models.} The NLP paper reports only accuracy, and the papers that measure $\lmax$ use small vision models.
\end{enumerate}

% ============================================================================
\section{Our contribution (planned)}
\label{sec:contribution}
% ============================================================================

\subsection{Goal}

Close the gap in point 3 above. We fine-tune a pre-trained transformer on GLUE with every optimizer covered in this talk, under one fair protocol.
For each run we measure \emph{both} generalization and curvature. This lets us ask Kaur et al.'s question in NLP.

\subsection{Research questions}

\begin{enumerate}[label=\textbf{RQ\arabic*.}, leftmargin=*]
  \item \textbf{Accuracy.} At an equal compute budget, do the SAM variants and Entropy-SGD beat AdamW when fine-tuning on small GLUE tasks?
        Do the gains grow when training data is scarce, as Bahri et al.\ found for T5?
  \item \textbf{Flatness.} Do these optimizers actually reach lower $\lmax$ and $\tr(H)$ on a transformer? Bahri et al.\ never checked.
  \item \textbf{Does flatness predict generalization in NLP?} Across all runs, how well does $\lmax$ (or the trace) rank the dev scores?
        Does Kaur et al.'s SAM result (benefit at small batch, none at large batch) also hold for transformers?
  \item \textbf{The X-SAM diagnostic.} Is the update also nearly orthogonal to the top Hessian eigenvector in transformers, as X-SAM found for CNNs?
\end{enumerate}

None of the papers we present runs Entropy-SGD, Eigen-SAM or X-SAM on transformers. RQ2--RQ4 have not been studied for NLP in any of them.

\subsection{Experimental design}

\begin{center}\small
\begin{tabularx}{\textwidth}{@{}p{3.2cm}L@{}}
\toprule
\textbf{Model} & BERT-base (uncased), fine-tuned \\
\textbf{Tasks} & RTE (2.5k train), MRPC (3.7k), CoLA (8.5k), and SST-2 subsampled to 5\% and 20\%. Bahri et al.\ report SAM's largest relative gains at these sampling rates (5\%: Tables 6--7; $\approx$20\%: Fig.~1), though on SuperGLUE and CBQA, not GLUE \\
\textbf{Base recipe} & AdamW, batch 32, 3 epochs, max length 128, 10\% warmup and linear decay, weight decay 0.01 (Devlin et al., Sec.~4.1) \\
\textbf{Optimizers} & AdamW (baseline); SAM, GSAM, WSAM, Eigen-SAM, X-SAM (all with AdamW as the base optimizer); Entropy-SGD in its Entropy-Adam form (Sec.~5.4.2: $L=5$, fixed $\gamma$, $\beta_1=0.5$) \\
\textbf{Tuning} & The protocol of the SAM-variant papers, with each method's grid from its own paper. (1) AdamW lr $\in\{2,3,4,5\}\times10^{-5}$ (Devlin et al.), then fixed for all methods (GSAM App.~B.1). (2) SAM $\rho\in\{0.02,0.05,0.1,0.15,0.2,0.3\}$ with the ascent on $1/4$ of the batch (Bahri et al.). (3) SAM's best $\rho$ reused by GSAM ($\alpha\in\{0.1,0.2,0.3\}$, $\rho$ tied to the lr), WSAM ($\gamma\in\{0.80,\dots,0.98\}$), Eigen-SAM ($\alpha\in\{0.05,0.1,0.2\}$, $p=100$, $q=5$) and X-SAM ($\alpha=0.2$). Selection on the dev set with one seed \\
\textbf{Fair comparison} & Equal number of gradient evaluations per run, as in Eigen-SAM (App.~E) and Entropy-SGD (``effective epochs'', Sec.~5.2): an update costs 1 (AdamW), 1.25 (SAM), 2 (GSAM, WSAM), 2.11 (Eigen-SAM, X-SAM) or $L=5$ (Entropy-Adam); 5 seeds of each method's best setting; mean $\pm$ std \\
\textbf{Generalization} & Standard dev metric per task (accuracy, F1 for MRPC, Matthews correlation for CoLA) at the final checkpoint, and the train--dev gap \\
\textbf{Curvature} & At the final checkpoint, on a fixed 1k-example training subset with dropout off: $\lmax$ by power iteration, $\tr(H)$ by Hutchinson, and SAM sharpness at fixed $\rho$. During training, at 5 points: $\lmax$ and the angle between $v_1$ and both the gradient (X-SAM's diagnostic) and the actual update \\
\textbf{Kaur-style test} & Kaur et al.'s Sec.~3.3 protocol on SST-2 (20\%): constant lr, train to 99\% train accuracy, $\rho\in\{0,0.02,\dots,0.1\}$, two learning rates, a small (32) and a large (256) batch, 4 seeds. Plus: scatter $\lmax$ against dev score over all runs and report rank correlation \\
\bottomrule
\end{tabularx}
\end{center}

\paragraph{Note on Adam.} AdamW rescales every coordinate, so the curvature it ``sees'' is the preconditioned Hessian, not $H$.
Where feasible we also report $\lmax$ of the preconditioned Hessian.

\paragraph{Unavoidable deviations from the source papers.} The model (BERT, not T5 or a CNN) and the tasks.
For Entropy-Adam, the paper's SGLD step $\eta'\in[0.1,1]$ is meant for SGD on CNNs and would destroy pre-trained weights, so we use Adam with a fixed $\eta'$ equal to the lr. Our $\gamma$ grid is $\{0.003,0.01,0.03\}$, and we record the paper's tuning signal $\|x-\mu\|/\|\nabla f\|$.
X-SAM does not state $p$, so we use Eigen-SAM's $p=100$.
We report the final checkpoint, where curvature is measured, rather than Bahri et al.'s best checkpoint.
The Kaur-style test uses AdamW and a large batch of 256 (2\% of the data, not 10\%), since larger batches would need gradient accumulation.

\subsection{Status}

\begin{itemize}
  \item \textbf{Done:} PyTorch implementations of SAM, GSAM, WSAM, Eigen-SAM, X-SAM and Entropy-SGD/Entropy-Adam behind one \texttt{optimizer.step(closure)} interface,
        with power iteration for $\lmax$; curvature tools ($\lmax$, Hessian trace) validated against the exact Hessian; GLUE fine-tuning scripts;
        a single runner (\texttt{Curvature-GLUE/run\_glue.py}) with compute matching, task-specific metrics, and final-checkpoint and in-training curvature measurements;
        a job queue (\texttt{sweep.py}) for the tuning stages above and for the Kaur et al.\ vision reproduction.
  \item \textbf{To do:} the runs and analysis; $\lmax$ of the preconditioned Hessian.
\end{itemize}

\subsection{What each outcome would mean}

\begin{itemize}
  \item \textbf{Flatter and better:} supports (H) for NLP fine-tuning and gives Bahri et al.'s gains a geometric explanation.
  \item \textbf{Better but not flatter, or $\lmax$ uncorrelated with dev score:} extends Kaur et al.'s conclusion to transformers. The gains then need another explanation.
  \item \textbf{Flatter but not better:} flatness-seeking alone does not transfer to fine-tuning; that is itself a useful negative result.
\end{itemize}

% ============================================================================
\appendix
\section{Datasets, tasks and metrics}
\label{app:data}
% ============================================================================

\begin{longtable}{@{}p{2.8cm}p{6.4cm}p{3.0cm}p{2.0cm}@{}}
\toprule
\textbf{Dataset} & \textbf{Task} & \textbf{Size} & \textbf{Used in} \\
\midrule
\endhead
MNIST & handwritten digits, $28{\times}28$ grayscale, 10 classes & 60k / 10k & I \\
Fashion-MNIST & same format, 10 clothing categories & 60k / 10k & II, III \\
CIFAR-10 & $32{\times}32$ colour photos, 10 classes & 50k / 10k & I, II, III \\
CIFAR-100 & same images, 100 classes & 50k / 10k & II \\
Penn Treebank & word-level language modelling & $\approx$930k / 74k / 82k words & I \\
\emph{War and Peace} & character-level language modelling & $\approx$3M characters & I \\
GLUE & 8 sentence(-pair) tasks: CoLA (grammaticality), SST-2 (sentiment), MRPC, QQP (paraphrase), STS-B (similarity), MNLI, QNLI, RTE (inference) & 2.5k--393k per task & II, ours \\
SuperGLUE & 8 harder tasks: BoolQ (yes/no QA), CB (entailment), COPA (causal reasoning), MultiRC, ReCoRD (reading comprehension), RTE, WiC (word sense), WSC (coreference) & 250--100k per task & II \\
Closed-book QA & Natural Questions, WebQuestions, TriviaQA: answer with no supporting passage & & II \\
TyDiQA-GoldP & extractive QA given the passage, 9 typologically diverse languages & & II \\
\bottomrule
\end{longtable}

\paragraph{Metrics.} \textbf{Accuracy / error rate.} \textbf{Perplexity} $= \exp(\text{mean negative log-likelihood per token})$; lower is better.
\textbf{Exact match (EM)} credits only identical answers, while \textbf{F1} credits overlapping tokens. \textbf{Matthews correlation} (CoLA) handles class imbalance;
\textbf{Pearson/Spearman} correlation is used for STS-B. GLUE and SuperGLUE scores are averages over tasks.

% ============================================================================
\section*{References}
\addcontentsline{toc}{section}{References}
% ============================================================================
\begin{enumerate}[label={[\arabic*]},leftmargin=*,itemsep=1pt]\small
  \item P.~Chaudhari et al. Entropy-SGD: Biasing gradient descent into wide valleys. \emph{ICLR}, 2017.
  \item P.~Foret, A.~Kleiner, H.~Mobahi, B.~Neyshabur. Sharpness-aware minimization for efficiently improving generalization. \emph{ICLR}, 2021.
  \item D.~Bahri, H.~Mobahi, Y.~Tay. Sharpness-aware minimization improves language model generalization. \emph{ACL}, 2022. arXiv:2110.08529.
  \item J.~Kwon, J.~Kim, H.~Park, I.~K.~Choi. ASAM: Adaptive sharpness-aware minimization for scale-invariant learning of deep neural networks. \emph{ICML}, 2021.
  \item J.~Zhuang et al. Surrogate gap minimization improves sharpness-aware training. \emph{ICLR}, 2022.
  \item Y.~Yue, J.~Jiang, Z.~Ye et al. Sharpness-aware minimization revisited: Weighted sharpness as a regularization term. \emph{KDD}, 2023.
  \item H.~Luo et al. Explicit eigenvalue regularization improves sharpness-aware minimization. arXiv:2501.12666, 2025.
  \item H.~Duan, Y.~Chen, L.~Guan. X-SAM: Boosting sharpness-aware minimization with dominant-eigenvector gradient correction. arXiv:2601.10251, 2026.
  \item S.~Kaur, J.~Cohen, Z.~C.~Lipton. On the maximum Hessian eigenvalue and generalization. arXiv:2206.10654, 2023.
  \item L.~Dinh, R.~Pascanu, S.~Bengio, Y.~Bengio. Sharp minima can generalize for deep nets. \emph{ICML}, 2017.
  \item J.~Cohen, S.~Kaur, Y.~Li, J.~Z.~Kolter, A.~Talwalkar. Gradient descent on neural networks typically occurs at the edge of stability. \emph{ICLR}, 2021.
  \item N.~S.~Keskar et al. On large-batch training for deep learning: Generalization gap and sharp minima. \emph{ICLR}, 2017.
  \item M.~Hardt, B.~Recht, Y.~Singer. Train faster, generalize better: Stability of stochastic gradient descent. \emph{ICML}, 2016.
  \item M.~Welling, Y.~W.~Teh. Bayesian learning via stochastic gradient Langevin dynamics. \emph{ICML}, 2011.
  \item S.~Hochreiter, J.~Schmidhuber. Flat minima. \emph{Neural Computation}, 1997.
  \item B.~A.~Pearlmutter. Fast exact multiplication by the Hessian. \emph{Neural Computation}, 1994.
\end{enumerate}

\end{document}
